\section{Literature Mapping Strategy}

The related-work analysis in this qualification document is guided by the
reference material stored in the local folder \texttt{artigos\_referencia}. The
folder contains papers organized into thematic groups, CSV files with article
classifications, keyword matrices, and automatically generated group
descriptions. These automatic groupings were produced from an
article-by-controlled-expression matrix weighted with TF-IDF and clustered with
AGNES. Therefore, they are useful as a navigation instrument, but they are not
treated as final scientific interpretation.

The search strategy was structured around four logical blocks executed across
five major digital libraries: IEEE Xplore, ACM Digital Library, Scopus,
SpringerLink, and Google Scholar.

\begin{description}
    \item[Block A (Source Code):] \texttt{"source code" OR "software code" OR "program code"}
    \item[Block B (Analysis):] \texttt{"code analysis" OR "software analysis" OR "program analysis"}
    \item[Block C (Techniques):] \texttt{"machine learning" OR "deep learning" OR "static analysis" OR "dynamic analysis"}
    \item[Block D (Visual Representation):] \texttt{"visual representation" OR "code visualization" OR "minimap"}
\end{description}

For queries targeting visual paradigms, the blocks were combined as
\textit{(A) AND (B) AND (C) AND (D)}. To cover broader analysis methodologies,
Block D was omitted in complementary search runs.

Reference management relied on Zotero for automated metadata harvesting. Study
selection and blinded screening were managed through the Rayyan platform.
The filtering followed three progressive stages: (i)~title and abstract
screening produced an initial set of 270 candidates; (ii)~a keyword-based
refinement guided by a control group of 7 seminal baseline papers, from which 39
structural index terms were extracted, narrowed the pool to 155 articles;
(iii)~negative filtering eliminated secondary studies and out-of-scope formats
(survey, review, mapping study, tutorial, editorial, poster, and works centered
on security or UX), reducing the active set to 96 articles. Full-text inspection
with iterative peer validation then produced the final corpus of 56 selected
studies.

The literature was organized around seven analytical dimensions:

\begin{enumerate}
    \item the software-engineering problem addressed by each work;
    \item the representation used for source code or software artifacts;
    \item the degree of language dependence;
    \item the learning or visualization method adopted;
    \item the dataset or empirical basis;
    \item the evaluation strategy;
    \item the limitations that motivate the minimap-based proposal.
\end{enumerate}

This organization is intentionally problem-oriented rather than chronological.
The goal is not only to list existing tools or models, but to identify where
source code minimaps can contribute. The following sections synthesize the main
families of work.

\section{Interpretation of the Local Article Groups}

The local article groups provide an initial map of the literature. The largest
group concentrates on software visualization, source code analysis, static
analysis, software quality, maintenance, and evolution. This group is important
because it establishes the software-engineering context in which minimaps should
be positioned. It includes works that visualize dependencies, architecture,
tests, source-code relations, microservice boundaries, fault localization, and
quality indicators. These works show that visual and structural abstractions are
useful, but they usually remain tied to specific analysis goals or human-facing
views.

A smaller but central group is directly related to minimaps, Local Binary
Patterns, feature detection, and visual signatures. This group contains the
initial minimap study and therefore functions as the seed of the doctoral
proposal. Its role in the literature review is not only to cite previous work
by the candidate, but to identify what remains open after that first study:
larger datasets, deep models, explainability, stronger anonymization analysis,
and tool integration.

Another group emphasizes computer vision, Vision Transformers, ResNet,
Transformer-based models, and source-code representations such as CV4Code. This
group helps connect minimap analysis to modern visual learning. It also
supports the methodological move from handcrafted texture descriptors to
end-to-end learning. The relevance of this group is not that source code
minimaps are natural images, but that computer vision methods can process
spatial patterns and learn hierarchical representations.

The final group contains works related to large language models, conversational
data analysis, human-AI interaction, code verification, and visual programming.
This group is more peripheral to the current qualification proposal, but it is
useful for future positioning. As software-engineering tools increasingly adopt
AI assistants, minimap-based analysis may become one component of a multimodal
toolchain that combines visual, textual, structural, and conversational feedback.

This interpretation of the groups leads to a practical decision for the thesis:
the related-work chapter should not be organized by the automatic group numbers.
Instead, it should translate those groups into conceptual families. The families
used in this chapter are software visualization, machine learning for code,
visual code representations, software quality, explainability, and tool support.

\section{Software Visualization and Visual Code Inspection}

Software visualization is the most direct conceptual predecessor of this
project. It assumes that visual abstractions can help developers and researchers
reason about software systems in ways that raw text alone cannot support. In
early and classical work, visualization often focused on dependencies,
architectural structures, metrics, or evolution. Pinzger et al.
\cite{pinzger2008tool} proposed visual support for understanding source code
dependencies, while SourceMiner Evolution supported feature evolution
comprehension through coordinated views \cite{novais2013sourceminer}. SArF Map
used map metaphors to visualize software architecture from feature and layer
viewpoints \cite{kobayashi2013sarf}.

These works are important because they show that visual representations can
externalize structural information that is difficult to grasp from files alone.
However, their main output is designed for human inspection. A developer or
architect examines the visualization and draws conclusions. In contrast, this
work treats minimaps as visual signatures that can be processed by
models. The human may still inspect the image or the explanation, but the first
analytical step is computational.

Recent visualization tools reinforce the relevance of this area. FineCodeAnalyzer
provides multi-perspective visualization for fine-grained source code analysis
\cite{qayum2022finecodeanalyzer}. SFLVis supports visual analysis of software
fault localization results \cite{sun2024sflvis}. Code-smell detection and
visualization have also been reviewed systematically, showing that visual support
remains useful for quality-related tasks \cite{reis2022code}. ScaMaha parses,
analyzes, and visualizes object-oriented systems \cite{almsiedeen2025scamaha},
and CodeSketch explores real-time multi-language visualization with UML and
animation support \cite{shrivastava2025codesketch}.

CodePanorama is particularly close to this project because it explores
language-agnostic visual code inspection \cite{etter2022codepanorama}. It
supports the intuition that zoomed-out visual representations can expose
patterns across source artifacts. Nevertheless, CodePanorama remains primarily a
visual inspection tool. The minimap approach proposed here extends this line by
turning compact visual representations into data for supervised learning,
explainability, and IDE-integrated inference.

\section{Machine Learning Representations for Source Code}

A second body of work studies how source code can be represented for machine
learning. Allamanis et al. \cite{allamanis2018survey} review the broader field
of machine learning for big code and identify several recurring tasks and
representations. Token-based and sequence-based models treat code as a language,
while tree-based and graph-based models capture syntactic and semantic structure.
Code2vec, for example, represents code snippets using paths in abstract syntax
trees \cite{alon2019code2vec}. CodeBERT and GraphCodeBERT use transformer-based
pretraining for code and natural language, with GraphCodeBERT incorporating data
flow information \cite{feng2020codebert,guo2021graphcodebert}.

These approaches have strong semantic and syntactic advantages. They can capture
identifier use, API calls, syntactic roles, data flow, and natural-language
comments. However, this expressiveness often comes with preprocessing
requirements. Models may depend on tokenizers, parsers, language-specific
grammars, valid syntax, or large pretraining corpora. Such requirements are
reasonable for many tasks, but they can complicate analysis across repositories
containing many languages and non-code artifacts.

Graph-based methods provide another important representation family. Nadim et
al. \cite{nadim2022sourcecodegraphs} use structural properties of source code
graphs for just-in-time bug prediction. Graphs are highly informative when the
task depends on dependencies, calls, control flow, or structural relations. Yet
graphs require extraction procedures and normalization decisions that may vary
across languages and tools. SparseCoder represents a more recent effort to make
source code analysis more efficient with sparse attention and learned token
pruning \cite{yang2025sparsecoder}. Such advances reduce computational cost, but
they remain within a textual or token-based modeling tradition.

The minimap proposal differs in its first representation layer. Instead of
parsing source code into language-aware structures, it transforms files into
images. This loses explicit semantics, but it gains uniformity across many file
types. The proposal should therefore be understood as complementary to
token-based, tree-based, and graph-based methods. A practical pipeline may use
minimaps to quickly classify and prioritize artifacts, then apply deeper static
or semantic analysis only where needed.

\section{Computer Vision and Visual Representations of Code}

The idea of using visual representations for code has gained momentum with work
such as CV4Code \cite{shi2023cv4code}. CV4Code investigates source code
understanding through visual code representations, showing that image-like
encodings can support learning tasks that are usually addressed with textual or
syntactic models. This work is important because it validates the broader
premise that source code contains visual signals useful for machine learning.

A significant finding from the broader visual-code literature is the transition
toward deep visual computing. Early approaches to visual code analysis relied on
handcrafted texture descriptors such as Local Binary Patterns. More recent work
demonstrates increasing adoption of end-to-end computer vision architectures:
Convolutional Neural Networks, residual learning backbones (ResNet), and,
more recently, Vision Transformers (ViTs) are reported in the literature for
processing anonymized or visual code representations directly. The patch-based
attention mechanism of ViTs aligns naturally with spatial regions of a minimap
such as indentation bands, delimiter zones, and block boundaries, making them a
candidate architecture for future experiments.

In this work, the large-scale study (Study 2) evaluated ResNet-18
and EfficientNet-B0. The MinimapSignatures backend uses a custom
MultiTaskMinimapCNN with a shared convolutional backbone and independent
classification heads per target. Transformer-based and hybrid architectures
remain a direction for future model comparison rather than a current
contribution.

The minimap approach shares this premise but emphasizes a different set of
research concerns. First, minimaps are designed to be compact and editor-like,
with fixed-size variants such as 128 by 128 pixels. Second, the approach includes
anonymization as a central step, not only as a preprocessing option. Third, it
examines several prediction objectives: Type, Project, and Author. Fourth, it
uses explainability to investigate the visual regions learned by the models.
Fifth, it integrates the representation into IDE plugins through
MinimapSignatures.

The methodological inspiration also comes from domains beyond software
engineering. In music classification, spectrograms transform audio signals into
visual representations that can be analyzed by texture descriptors and neural
models \cite{costa2011music,costa2012music,costa2017cnn}. Similar strategies
have been used for sound-image classification in other domains
\cite{nanni2018bird,pereira2019representation}. These works suggest that a
domain need not be inherently visual to benefit from visual machine learning.
The key requirement is that the visual transformation preserves task-relevant
patterns.

For source code minimaps, the relevant patterns are not semantic objects but
layout regularities, symbol distributions, indentation structures, and repeated
regions. This makes the representation closer to texture and document-image
analysis than to natural-object recognition. Consequently, both handcrafted
texture descriptors and convolutional neural networks are relevant baselines.

\section{Software Quality, Code Smells, and Architecture-Specific Patterns}

One of the long-term goals of this work is to extend minimap-based
analysis toward software quality indicators, code smells, and architectural
deviations. The literature on code smells is therefore relevant. Aniche et al.
\cite{aniche2018code} show that Model--View--Controller architectures have
specific code smells, highlighting that quality problems often depend on
architectural roles and framework conventions. A generic rule may miss problems
that become visible only when the expected role of an artifact is known.

This observation is important for minimap research. If files with different
roles produce different visual signatures, then minimaps may help identify
artifacts whose structure deviates from project or architecture expectations.
For example, a controller, repository, entity, test file, or configuration file
may have visual regularities that are not sufficient to prove correctness but
are useful as weak signals. A minimap-based model could flag unusual files for
human review or for deeper static analysis.

However, quality assessment is more demanding than Type classification. A file
can visually resemble a valid controller and still be poorly designed. A
minimap can reveal structure, density, repetition, and layout, but it cannot
directly infer all semantic properties. Therefore, quality-related extensions
must be designed carefully, possibly combining minimap signals with metrics,
static-analysis outputs, or repository metadata.

\section{Explainability for Visual and Software Models}

Explainability is relevant in two ways. First, deep learning models applied to
minimaps must be inspected to determine whether they use plausible evidence.
Second, if minimap-based tools are to provide developer-facing feedback, the
feedback should not be limited to opaque labels and confidence scores.

Integrated Gradients provides one method for attributing predictions to input
features \cite{integratedgradients2017}. Broader XAI surveys discuss the
strengths and limitations of attribution, surrogate models, and local
explanations \cite{adadi2018xai}. In minimap analysis, attribution maps can
show whether predictions are influenced by indentation zones, headers,
boilerplate, dense generated regions, or file endings. This is valuable for
interpretation, but it must not be overstated. Attribution maps are not proofs
that a model understands source code; they are evidence about sensitivity and
salient visual regions.

The large-scale minimap study uses Integrated Gradients to compare explanations
across Type, Project, and Author objectives \cite{gebara2025ciarp}. Early
findings suggest that Type predictions often rely on structural bands,
Project predictions may emphasize repository conventions or boilerplate, and
Author predictions can be more diffuse. These results support the idea that
different objectives depend on different visual cues, but they also raise
important methodological questions about leakage, convention, and over-reliance
on non-semantic artifacts.

\section{Tool Support and IDE Integration}

A recurring limitation of research prototypes in software engineering is that
they remain detached from everyday development environments. Developers do not
usually analyze code through separate experimental pipelines; they work in IDEs,
version-control platforms, code-review systems, and continuous-integration
tools. For a minimap-based approach to be practically relevant, it must be
operationalized in such environments.

The tool literature includes many examples of visualization and analysis tools,
but fewer works connect compact visual representations, trained models, and
multi-IDE workflows. MinimapSignatures addresses this gap by using Visual Studio
Code and IntelliJ IDEA plugins to generate anonymized minimaps locally and send
them to a backend for prediction \cite{gebara2026minimapsignatures}. This design
separates IDE-specific interaction from model inference, allowing models to
evolve independently from the plugins.

The preliminary evaluation of MinimapSignatures is functional. It validates that
the complete workflow can be executed: source capture, minimap generation,
anonymization, HTTP communication, backend inference, and IDE feedback. This is
an important first step, but it is not yet a user study. Future work must
evaluate usefulness, trust, cognitive effort, and impact on development tasks.

\section{Application Domains Identified in the Mapping}

The systematic mapping reveals that source code minimaps have expanded well
beyond their historical role as IDE navigation widgets. The surveyed studies show
a strong concentration of minimap applications across three principal domains:

\begin{description}
    \item[Aesthetic and Structural Categorization.] Identifying functional
    stereotypes and geometric structural rhythms within multi-language
    repositories. Minimap classifiers can distinguish file roles based on visual
    patterns of indentation, delimiter density, and block geometry without
    exposing lexical content.

    \item[Author and Provenance Attribution.] Utilizing stylistic visual
    fingerprinting—such as individual indentation habits and spacing
    patterns—to map code authorship without exposing raw lexical tokens. This
    objective is treated as a research proxy for style and convention analysis.

    \item[Project and Ecosystem Clustering.] Discerning architectural similarities
    and matching files to specific project families or software components based
    on layout regularities and visual conventions unique to each codebase.
\end{description}

These domains confirm that minimap representations encode information at
multiple granularities: file type, project identity, and contributor style.
They also motivate the three prediction objectives (Type, Project, Author)
adopted in this work.

\section{Synthesis of Gaps}

Table~\ref{tab:related-synthesis} summarizes how the main related-work families
connect to this project and which gaps remain.

\begin{table}[ht]
\centering
\caption{Synthesis of related-work families and research gaps.}
\label{tab:related-synthesis}
\begin{tabular}{p{0.22\textwidth}p{0.33\textwidth}p{0.32\textwidth}}
\toprule
\textbf{Family} & \textbf{Typical contribution} & \textbf{Gap addressed by this project} \\
\midrule
Software visualization & Human-facing views for dependencies, architecture, evolution, or faults. & Use compact visual artifacts as machine-readable signatures. \\
Machine learning for code & Token, tree, graph, or transformer representations. & Provide a language-light visual representation that can include non-code artifacts. \\
Visual code representations & Treat code as image-like data. & Add anonymization, scale, explainability, and IDE integration around minimaps. \\
Quality and smells & Detect or visualize maintainability problems. & Investigate minimap signals as weak evidence for quality and architectural deviation. \\
Explainable AI & Attribute model predictions to inputs. & Interpret what visual regions of minimaps support software-analysis predictions. \\
IDE tool support & Bring analysis closer to developers. & Operationalize minimap generation and inference in multi-IDE workflows. \\
\bottomrule
\end{tabular}
\fautor
\end{table}

From this synthesis, the core research opportunity becomes clear. Source code
minimaps may provide a compact bridge between software visualization and
machine learning for code. They are visual enough to capture spatial and
structural patterns, simple enough to generate across heterogeneous repositories,
and compatible with both handcrafted descriptors and deep visual models. The
remaining challenge is to evaluate them rigorously, interpret their behavior,
understand their limitations, and integrate them into practical workflows.

\section{Positioning of the Proposed Research}

The proposed research can be positioned through three contrasts.

First, compared with traditional software visualization, minimap-based analysis
changes the primary consumer of the representation. A dependency graph, treemap,
or architecture map is usually designed so that a human can inspect it directly.
A source code minimap may be inspected by a human, but its first role is to
serve as data for a model. This distinction affects the design criteria. A human
visualization should be expressive, readable, and interactive. A minimap
signature should be compact, reproducible, and discriminative.

Second, compared with token, tree, graph, and transformer-based code models,
minimap-based analysis changes the first abstraction layer. Instead of
constructing a language-aware representation, it renders source artifacts into a
visual space. This makes the representation weaker semantically but broader
operationally. It can include code, markup, configuration, scripts, templates,
and generated artifacts under the same rendering procedure. This is particularly
useful for repositories whose maintenance problems span many kinds of files.

Third, compared with standalone machine learning pipelines, MinimapSignatures
connects the representation to developer environments. This is important because
software-engineering research often stops at the point where a model reports a
metric. The tool line asks what happens next: how a developer obtains the
prediction, how the result is shown, which files can be analyzed, and how a
backend can evolve independently from IDE clients.

These contrasts define the distinctive contribution of the thesis. It is not
only a new classifier, a new visualization, or a new plugin. It is a pipeline
that links visual representation, anonymization, empirical evaluation,
explainability, and tool support.

\section{Open Problems Derived from the Literature}

The literature also reveals open problems that the thesis must address. Three of
these gaps recur consistently across the surveyed studies and can be named
precisely:

\paragraph{Tooling Detachment.}
Visual source code machine learning pipelines remain almost entirely
disconnected from real-time developer workflows. There is a lack of automated
plugin infrastructures capable of capturing, anonymizing, and transmitting active
editing sequences directly from live modern IDEs. Offline accuracy metrics do not
capture the operational friction of live software engineering, where code
structures are constantly incomplete or undergoing syntax updates.

\paragraph{Absence of User-Centered Evaluation.}
Empirical evidence regarding usability, cognitive load, developer trust in AI
recommendations, and practitioner interaction paradigms is exceptionally scarce.
Almost no studies measure how real developers interact with visual AI feedback.
The community cannot yet assess whether minimap-based tools are useful,
interpretable, or trustworthy from a practitioner perspective.

\paragraph{The Single-File Scale Barrier.}
Current prototypes operate almost exclusively on isolated, single-file
representations at a fixed visual scale. Mechanisms supporting multi-file
modules, project-level collection analysis, adaptive multi-scale architectures,
and user-driven model retraining are virtually absent in the literature. This
limit prevents minimap analysis from addressing tasks that inherently span
multiple files or directory trees.

Beyond these three structural gaps, the thesis must also address:

\begin{itemize}
    \item \textit{Representation fidelity}: minimaps preserve layout and symbols but
    do not explicitly encode type hierarchies, API usage, or data flow. The thesis
    must identify which tasks are appropriate for minimaps and which require
    complementary representations.

    \item \textit{Privacy and leakage}: anonymization reduces exposure but is not
    a formal privacy guarantee. Generated files, framework templates, or
    distinctive repository conventions may remain identifiable. Stronger
    anonymization modes must be evaluated.

    \item \textit{Generalization}: strong performance inside one corpus does not
    guarantee transfer to other ecosystems. Repository-aware splits, ablations,
    and per-class analysis are necessary.
\end{itemize}

These problems shape the research plan presented later in this document.
